(1) 取 $|+z\rangle=(1,0)^T$，$S_x=(\hbar/2)\begin{pmatrix}0&1\\1&0\end{pmatrix}$，其本征态 $|\pm x\rangle=(|+z\rangle\pm|-z\rangle)/\sqrt2$。投影概率 $P(\pm x|+z)=|\langle\pm x|+z\rangle|^2=1/2$，两路相加为一。
(2) 在首次得到 $+x$ 的条件下状态经投影变为 $|+x\rangle=(|+z\rangle+|-z\rangle)/\sqrt2$。其后测 $S_z$ 的概率分别为 $P(+z|+x)=P(-z|+x)=1/2$，与原来的确定 $+z$ 不同，说明不对易测量会改变后续统计。
(3) 若前次 $\pm x$ 结果均保留而不记录，第一次投影后状态为混合态 $\rho=\frac12|+x\rangle\langle+x|+\frac12|-x\rangle\langle-x|=I/2$。随后测 $+z$ 的概率为 $\operatorname{Tr}[\rho|+z\rangle\langle+z|]=1/2$，也可用全概率公式 $\frac12\times\frac12+\frac12\times\frac12=1/2$。
(4) 沿单位矢量 $\hat n$ 的正自旋纯态密度矩阵为 $\rho_n=(I+\hat n\cdot\boldsymbol\sigma)/2$，沿 $\hat m$ 的正结果投影为 $P_m=(I+\hat m\cdot\boldsymbol\sigma)/2$。所以条件概率 $\operatorname{Tr}(\rho_nP_m)=(1+\hat n\cdot\hat m)/2=\cos^2(\theta/2)$，当两轴相同为一、相反为零、正交为二分之一。
